% Section 7.4, from lectures/L07/appendix/c_exercises.md. The worked solutions are not
% printed; the aside says where they are.

\section{Exercises}
\label{c7:sec:exercises}

Eight, ending with the Cross-check. That one settles which resistance the emitter factor multiplies, and the answer differs from the obvious one by a factor of ten.

\begin{aside}
\textbf{How to check your work.} A \kindtag{Code} exercise is judged by the suite that ships with it: run \code{make test} at the root of the companion repository, with \code{AEL_DIR} pointing at your toolkit. A suite you have not started is reported as \code{SKIP} rather than as a failure, so the command is worth running from the first day. The hand calculations and designs are checked against the toolkit functions each one names.

\textbf{The solutions are not in this book.} They are published in full in the companion repository, in \file{lectures/L07/appendix}, including the plausible wrong answer where there is one. Read them once you have your own answers, including the ones you are unsure about, because being unsure and right is a different thing from being unsure and wrong and the solutions say which is which.
\end{aside}

% ------------------------------------------------------------------
\recall{What linearising discards}
\label{c7:ex:1}
\begin{enumerate}
  \item State $r_e$, and say whether it is a resistor.
  \item How large may a base signal be before the straight-line model is misleading? Give a number and the reason for it.
  \item Name three things a small-signal model cannot describe.
  \item A stage is switched off. What happens to $r_e$?
\end{enumerate}

% ------------------------------------------------------------------
\recall{The four rules}
\label{c7:ex:2}
\begin{enumerate}
  \item List the four steps that turn a stage into its small-signal schematic.
  \item Why does a supply rail become a ground?
  \item A bias divider of 33 kilohm over 6.8 kilohm sits on the base of a stage whose base looks like 13 kilohm. Which dominates the input resistance, and by how much?
  \item Bypassing the emitter resistor recovers the gain. What else does it recover, and what does it throw away?
\end{enumerate}

% ------------------------------------------------------------------
\handcalc{The three results}
\label{c7:ex:3}
A stage at 1 mA with $R_C = 10$ kilohm, $R_E = 234$ ohm, $\beta = 50$.

\begin{enumerate}
  \item $r_e$, and the emitter factor.
  \item The gain, and what it would be with the emitter resistor bypassed.
  \item The input resistance looking into the base.
  \item The same three at 10 mA with the same resistors. Which of them changed, and which did not?
\end{enumerate}
\checkyourself{\code{intrinsicEmitterResistance}, \code{emitterFactor}, \code{gain}, \code{inputResistance}}

% ------------------------------------------------------------------
\handcalc{Two nodes, two answers}
\label{c7:ex:4}
For the same stage:

\begin{enumerate}
  \item The resistance looking into the collector, with $R_C$ removed.
  \item The stage's output resistance, with $R_C$ in place.
  \item The same two with the emitter resistor removed.
  \item State, in one sentence each, what degeneration did to each of the two quantities.
\end{enumerate}
\checkyourself{\code{resistanceIntoCollector}, \code{outputResistance}}

% ------------------------------------------------------------------
\design{A stage with a stated gain and bandwidth}
\label{c7:ex:5}
Design a common-emitter stage: 12 V supply, gain of 40 with the emitter resistor unbypassed, and an input corner above 3 MHz when driven from 600 ohm. Take $C_{bc} = 4$ pF.

\begin{enumerate}
  \item Choose a collector current and a collector resistor.
  \item Choose the emitter resistor for the gain.
  \item Compute the Miller capacitance and the input corner. Does it meet the requirement?
  \item If it does not, say what you would change, and give the two options with their costs.
\end{enumerate}
\checkyourself{\code{gain}, \code{millerCapacitance}, \code{intrinsicEmitterResistance}}

% ------------------------------------------------------------------
\codeexercise{The small-signal model}
\label{c7:ex:6}
Implement \code{ael::ssm} to the specification in \secref{c7:sec:b7}.

\textbf{\code{cascodeOutputResistance} must call \code{resistanceIntoCollector}} with $R_E = r_o$, rather than implementing a separate formula. The chapter claims a cascode is degeneration by $r_o$; writing it that way makes the code demonstrate the claim instead of restating it.

% ------------------------------------------------------------------
\codeexercise{The Early effect}
\label{c7:ex:7}
Multiply the forward transport current by $(1 + V_{CE}/V_A)$ in \code{ael::device::bjt}.

Then check it: solve a stage at two collector voltages a volt apart and confirm the collector current changes by about one part in a hundred, which is $1/V_A$ per volt. If the current does not move at all, $r_o$ is infinite and the Cross-check below cannot work.

% ------------------------------------------------------------------
\crosscheck{What the emitter factor multiplies}
\label{c7:ex:8}
A stage at 1 mA with $R_C = 10$ kilohm and $R_E = 234$ ohm, so $EF = 10$. Find its \textbf{output resistance}, four ways.

\begin{enumerate}
  \item \textbf{By the tempting rule.} $R_{out} = R_C \cdot EF$.
  \item \textbf{By the corrected closed form.} $R_C$ in parallel with the resistance looking into the collector.
  \item \textbf{By your solver.} Perturb the collector node by a small voltage with the input held, and divide the voltage change by the current change.
  \item \textbf{Then replace the collector resistor with a current mirror} and repeat legs 2 and 3.
\end{enumerate}

\task{What to expect}
\textbf{Leg 1 gives 100 kilohm. Legs 2 and 3 give about 9.9 kilohm.} A factor of ten, and it is the largest disagreement in this book.

\textbf{Leg 1 cannot be right, and you can see it without computing anything.} The output resistance is $R_C$ in parallel with something. A parallel combination is smaller than either part. So the answer cannot exceed 10 kilohm, and 100 kilohm does.

\textbf{Compare against the same stage with no emitter resistor: 9.09 kilohm.} So all that degeneration bought, at the cost of a factor of ten in gain, was \textbf{9 per cent} of output resistance. The boost is real, it is a factor of ten, and it is happening at a node where $R_C$ hides it.

\textbf{Leg 4 is where it stops being hidden.} With a current mirror the load is $r_o$, 100 kilohm rather than 10, and the numbers become 50 kilohm without degeneration and \textbf{89.6 kilohm} with it. The emitter factor is finally visible, because there is no longer a resistor swamping it.

\textbf{And that is the point of the correction.} It is not only that leg 1's number is wrong; it is that getting the node right explains why mirror loads exist at all. Without it, a mirror load is a technique that appears from nowhere. With it, it is the only way to collect something you have already paid for.

\textbf{If leg 3 gives exactly 10 kilohm}, the Early effect is missing from your device model, so $r_o$ is infinite and the transistor is a perfect current source. Check Exercise~\ref{c7:ex:7}.

\textbf{If leg 3 gives 100 kilohm}, the collector resistor is not in the netlist.

\textbf{If legs 2 and 3 differ by a few per cent}, that is expected: leg 2's closed form uses $R_E \parallel r_\pi$ where the solver uses the full nonlinear device, and the two agree to about the accuracy of $\beta$.
